\documentclass[11pt,a4paper]{article}

\usepackage[utf8]{inputenc}
\usepackage[T1]{fontenc}
\usepackage[spanish]{babel}
\usepackage[margin=2.2cm]{geometry}
\usepackage{xcolor}
\usepackage{tikz}
\usepackage{booktabs}
\usepackage{hyperref}
\usepackage{enumitem}
\usepackage{fancyhdr}
\usetikzlibrary{shapes.geometric, arrows.meta, positioning, fit, backgrounds}

% --- Brand-ish palette ---
\definecolor{brandprimary}{HTML}{6B4FA8}
\definecolor{brandaccent}{HTML}{A68FD6}
\definecolor{brandsoft}{HTML}{EEE7FA}
\definecolor{brandok}{HTML}{2E8B57}
\definecolor{brandwarn}{HTML}{C27500}

\hypersetup{
  colorlinks=true,
  linkcolor=brandprimary,
  urlcolor=brandprimary
}

\pagestyle{fancy}
\fancyhf{}
\fancyhead[L]{\textcolor{brandprimary}{\textbf{Astar Katar}}}
\fancyhead[R]{Programa Activaciones Diarias}
\fancyfoot[C]{\thepage}
\renewcommand{\headrulewidth}{0.4pt}

% --- TikZ node styles ---
\tikzset{
  startstop/.style   = {rectangle, rounded corners, minimum width=3.2cm, minimum height=0.9cm,
                        text centered, draw=brandprimary, fill=brandsoft, font=\small\bfseries},
  akbox/.style       = {rectangle, rounded corners=3pt, minimum width=3.6cm, minimum height=0.9cm,
                        text centered, text width=3.6cm, draw=brandprimary, fill=white, font=\small},
  decision/.style    = {diamond, aspect=2, minimum width=3.4cm, minimum height=0.6cm,
                        text centered, text width=3.2cm, draw=brandaccent, fill=brandsoft, font=\scriptsize},
  db/.style          = {cylinder, shape border rotate=90, aspect=0.3, minimum width=2.6cm, minimum height=1cm,
                        text centered, draw=brandprimary, fill=brandsoft, font=\small\bfseries},
  external/.style    = {rectangle, minimum width=3.2cm, minimum height=0.9cm, text centered,
                        draw=brandwarn, fill=white, font=\small\bfseries, dashed},
  arrow/.style       = {-{Stealth[length=2.5mm]}, thick, draw=brandprimary}
}

\begin{document}

% =================================================================
\begin{titlepage}
\centering
\vspace*{3cm}
{\Huge\bfseries\color{brandprimary} Programa Activaciones Diarias}\\[0.4cm]
{\Large Cómo funciona el sistema}\\[1.5cm]

\begin{tikzpicture}
  \fill[brandprimary] (0,0) circle (0.08);
  \fill[brandaccent] (0.6,0) circle (0.08);
  \fill[brandsoft] (1.2,0) circle (0.08);
\end{tikzpicture}

\vfill
{\large Documento explicativo para el dueño}\\[0.4cm]
{\small Astar Katar $\cdot$ \today}
\end{titlepage}

% =================================================================
\section{Resumen en una página}

El programa \textbf{Activaciones Diarias} es una suscripción de tiempo limitado
que da al cliente acceso a una biblioteca privada de audios guiados. No se vende
por pieza: el cliente paga un plan (1 día, 1 mes, 3 meses o 6 meses) y mientras
esté vigente puede escuchar \emph{toda} la biblioteca en línea.

\begin{center}
\begin{tikzpicture}[node distance=1.5cm]
  \node[startstop] (user) {Cliente};
  \node[akbox, right=of user] (pay) {Paga plan en Clip};
  \node[akbox, right=of pay] (access) {Escucha biblioteca};
  \node[db, right=of access] (exp) {Vence};

  \draw[arrow] (user) -- (pay);
  \draw[arrow] (pay) -- (access);
  \draw[arrow] (access) -- (exp);
  \draw[arrow, dashed] (exp.north) to[bend left=35] node[above, font=\scriptsize] {renueva} (pay.north);
\end{tikzpicture}
\end{center}

\subsection*{Puntos clave}
\begin{itemize}[leftmargin=*]
  \item El cobro se hace con \textbf{Clip}. Nosotros no guardamos datos de tarjeta.
  \item La suscripción se \textbf{activa sola} cuando Clip confirma el pago (webhook).
  \item Si el cliente renueva antes de vencer, \textbf{se suman los días} (no se pierde tiempo).
  \item Los audios \textbf{no se pueden descargar}: cada reproducción usa un enlace firmado que expira a los 10 minutos.
\end{itemize}

\newpage
% =================================================================
\section{Planes y precios}

Los cuatro planes viven en la base de datos (tabla \texttt{subscription\_plans}).
Modificarlos no requiere cambiar código.

\begin{center}
\begin{tabular}{@{}llrr@{}}
\toprule
\textbf{SKU} & \textbf{Nombre} & \textbf{Duración} & \textbf{Precio (MXN)} \\
\midrule
ACT-1D & Prueba 1 día & 1 día    & \$90 \\
ACT-1M & 1 mes        & 30 días  & \$390 \\
ACT-3M & 3 meses      & 90 días  & \$990 \\
ACT-6M & 6 meses      & 180 días & \$1{,}690 \\
\bottomrule
\end{tabular}
\end{center}

\noindent\textbf{Nota:} el plan de \emph{prueba 1 día} sirve para que el cliente
conozca la biblioteca antes de comprometerse con un plan largo.

% =================================================================
\section{Flujo 1: Compra de una suscripción}

Desde que el cliente elige un plan hasta que queda activo en su cuenta.

\begin{center}
\begin{tikzpicture}[node distance=1.1cm and 1.3cm]
  \node[startstop] (client) {Cliente elige plan};
  \node[akbox, below=of client] (web) {Sitio web crea referencia\\ \textit{SUB-ACT-\{uid\}-\{sku\}}};
  \node[external, below=of web] (clip) {Clip (pasarela)};
  \node[decision, below=of clip] (ok) {¿Pago exitoso?};
  \node[akbox, below left=1cm and 0.3cm of ok] (fail) {Cliente ve error\\ nada cambia en BD};
  \node[akbox, below right=1cm and 0.3cm of ok] (wh) {Webhook \texttt{transaction.succeeded}};
  \node[akbox, below=1.2cm of wh] (verify) {Verifica firma HMAC\\ (secreto compartido con Clip)};
  \node[decision, below=of verify] (dup) {¿Ya procesado?};
  \node[akbox, left=of dup] (skip) {Ignora duplicado};
  \node[db, below=of dup] (db) {Tabla \texttt{subscriptions}};
  \node[akbox, below=of db] (done) {Cliente ve acceso\\ en Mis Compras};

  \draw[arrow] (client) -- (web);
  \draw[arrow] (web) -- (clip);
  \draw[arrow] (clip) -- (ok);
  \draw[arrow] (ok) -| node[above, pos=0.25, font=\scriptsize] {no} (fail);
  \draw[arrow] (ok) -| node[above, pos=0.25, font=\scriptsize] {sí} (wh);
  \draw[arrow] (wh) -- (verify);
  \draw[arrow] (verify) -- (dup);
  \draw[arrow] (dup) -- node[above, font=\scriptsize] {sí} (skip);
  \draw[arrow] (dup) -- node[right, font=\scriptsize] {no} (db);
  \draw[arrow] (db) -- (done);
\end{tikzpicture}
\end{center}

\subsection*{¿Qué pasa si el cliente ya tiene una suscripción activa?}

No se crea una segunda. El sistema \textbf{extiende} la existente: suma los días
del nuevo plan a la fecha de vencimiento actual. La regla es
\emph{``nunca encoger''}: si el plan vence el 15 de junio y hoy renueva por 30 días,
el nuevo vencimiento es el 15 de julio (no desde hoy).

\newpage
% =================================================================
\section{Flujo 2: Escuchar los audios}

El cliente entra a \texttt{/activaciones-diarias/biblioteca}. El navegador
pide al servidor los enlaces. El servidor valida suscripción y firma cada URL
con una clave secreta. Los audios viven en un directorio privado; nunca se
sirven directo.

\begin{center}
\begin{tikzpicture}[node distance=1.1cm]
  \node[startstop] (a) {Cliente abre Biblioteca};
  \node[akbox, below=of a] (b) {Navegador pide lista firmada};
  \node[decision, below=of b] (c) {¿Suscripción activa?};
  \node[akbox, left=of c] (deny) {403 Forbidden\\ sin acceso};
  \node[akbox, below=of c] (d) {Servidor firma cada audio\\ (HMAC + exp 10 min)};
  \node[akbox, below=of d] (e) {Cliente reproduce audio};
  \node[decision, below=of e] (f) {¿Firma válida\\ y vigente?};
  \node[db, below=of f, align=center] (g) {Carpeta privada};
  \node[akbox, left=of f] (g2) {403 Forbidden};

  \draw[arrow] (a) -- (b);
  \draw[arrow] (b) -- (c);
  \draw[arrow] (c) -- node[above, font=\scriptsize] {no} (deny);
  \draw[arrow] (c) -- node[right, font=\scriptsize] {sí} (d);
  \draw[arrow] (d) -- (e);
  \draw[arrow] (e) -- (f);
  \draw[arrow] (f) -- node[above, font=\scriptsize] {no} (g2);
  \draw[arrow] (f) -- node[right, font=\scriptsize] {sí} (g);
\end{tikzpicture}
\end{center}

\subsection*{¿Por qué tanto cuidado?}

Sin firma, cualquier persona con el enlace podría compartir un audio o descargarlo.
Con firma:

\begin{itemize}[leftmargin=*]
  \item Cada enlace es \textbf{único por cliente} y \textbf{expira a los 10 minutos}.
  \item Compartir el enlace no sirve: expira antes de que el otro lo use.
  \item El reproductor tiene el botón de \textbf{descargar oculto} (no es el candado, pero desalienta al usuario casual).
  \item Al servidor solo llegan pedidos con firma válida \emph{y} suscripción vigente.
\end{itemize}

\noindent\textbf{Aviso honesto:} un técnico con herramientas de desarrollador
puede capturar el audio mientras suena. No existe protección perfecta en web.
El sistema disuade la copia casual; no imita DRM de Spotify.

\newpage
% =================================================================
\section{Flujo 3: Vista del cliente en ``Mis Compras''}

\begin{center}
\begin{tikzpicture}[node distance=1.1cm]
  \node[startstop] (l) {Cliente inicia sesión};
  \node[akbox, below=of l] (f1) {Sistema consulta 2 suscripciones};
  \node[akbox, below=of f1] (f2) {meditaciones + activaciones\_diarias};
  \node[decision, below=of f2] (d1) {¿Activaciones\\activa?};
  \node[akbox, left=of d1] (hide) {No se muestra banner};
  \node[akbox, below=of d1] (show) {Banner ``Activaciones Diarias''\\ con días restantes};
  \node[akbox, below=of show] (go) {Botón ``Ir a Biblioteca''};

  \draw[arrow] (l) -- (f1);
  \draw[arrow] (f1) -- (f2);
  \draw[arrow] (f2) -- (d1);
  \draw[arrow] (d1) -- node[above, font=\scriptsize] {no} (hide);
  \draw[arrow] (d1) -- node[right, font=\scriptsize] {sí} (show);
  \draw[arrow] (show) -- (go);
\end{tikzpicture}
\end{center}

% =================================================================
\section{Ciclo de vida de una suscripción}

\begin{center}
\begin{tikzpicture}[node distance=1.6cm]
  \node[startstop, fill=brandsoft] (nosub) {Sin suscripción};
  \node[startstop, right=of nosub, fill=white, draw=brandok, text=brandok] (act) {Activa};
  \node[startstop, right=of act, fill=white] (exp) {Vencida};

  \draw[arrow] (nosub) -- node[above, font=\scriptsize] {paga plan} (act);
  \draw[arrow] (act.north) to[bend left=25] node[above, font=\scriptsize] {renovación (suma días)} (act.north);
  \draw[arrow] (act) -- node[above, font=\scriptsize] {llega \texttt{end\_date}} (exp);
  \draw[arrow, dashed] (exp.south) to[bend left=35] node[below, font=\scriptsize] {vuelve a pagar} (act.south);
\end{tikzpicture}
\end{center}

\noindent Estado en la base de datos:
\begin{itemize}[leftmargin=*]
  \item \texttt{status = 'Activa'} + \texttt{end\_date > NOW()}: el cliente tiene acceso.
  \item \texttt{end\_date < NOW()}: vencida, se le quita el acceso al siguiente pedido.
  \item Se mantiene el historial completo en \texttt{subscription\_events} (auditoría de pagos).
\end{itemize}

% =================================================================
\section{Qué pasa si algo falla}

\begin{center}
\begin{tabular}{@{}p{5.8cm}p{9cm}@{}}
\toprule
\textbf{Síntoma} & \textbf{Qué revisar} \\
\midrule
El cliente pagó pero no ve acceso &
El webhook de Clip no llegó o la firma fue rechazada. Revisar logs del endpoint \texttt{/api/webhooks/clip}. \\
\addlinespace
El cliente pagó dos veces &
La tabla \texttt{subscription\_events} tiene la referencia de pago como única, así que el segundo webhook se ignora. Sí genera cargo en Clip: reembolsar desde el panel de Clip. \\
\addlinespace
El audio no se reproduce &
Verificar que el archivo existe en \texttt{private-audio/} y que la fila en \texttt{activaciones} apunta al \texttt{audio\_filename} correcto. \\
\addlinespace
La suscripción dice activa pero no deja escuchar &
Probablemente \texttt{end\_date} ya pasó pero \texttt{status} sigue en `Activa'. El endpoint de audio valida ambos. \\
\bottomrule
\end{tabular}
\end{center}

% =================================================================
\vfill
\begin{center}
\textit{Documento de referencia. Ante cualquier duda operativa, contactar al equipo técnico.}
\end{center}

\end{document}
